\section{Methodology}
\label{sec:methodology}

This section describes the pipeline that produced every result in Section~\ref{sec:results}: the datasets, the four encoder families used, the register mechanism, and the training procedure.

\subsection{Datasets and Preprocessing}
\label{sec:datasets}

The primary dataset is IAM~\cite{marti2002iam} at the line level, using the Aachen split~\cite{retsinas2022htrbestpractices}: $6{,}482$ training lines, $976$ validation lines, and $2{,}915$ test lines with no writer overlap between train and test. The character vocabulary is $79$ classes plus one CTC blank, giving $80$ output classes.

The generalization test set is READ2016 --- historical German handwriting distributed as $8{,}349/1{,}040/1{,}138$ train/val/test lines with an $88$-class vocabulary. Its train set is larger than IAM's, its character set is larger, and its script (historical German) differs from IAM's modern English; together this makes it a useful cross-dataset check on any conclusion reached on IAM. Ground-truth images are provided as RGB JPEGs; they are converted to 8-bit grayscale PNGs before use, to avoid a colour-conversion bug in the standard $1 - I/255$ inversion when applied to RGB input.

A synthetic corpus of roughly $950{,}000$ rendered lines is used only for the synthetic-pretraining lane in Section~\ref{sec:res_synth}: text sampled from CC100 is rendered with $555$ handwritten-style fonts using the pipeline of Nikolaidou~\etal~\cite{nikolaidou2023wordstylist}. It is not central to the register-token question.

All line images are preprocessed to $128 \times 1024$ grayscale tensors following the recipe of Retsinas~\etal~\cite{retsinas2022htrbestpractices}: aspect-ratio-preserving resize, symmetric white padding, per-pixel intensity inversion ($\hat I = 1 - I/255$), and per-image standardization. Training-time augmentation follows a moderate profile (small affine transforms, elastic distortion at $p{=}0.5$, mild morphological operations, brightness/contrast jitter) which the augmentation ablation in Section~\ref{sec:res_abl_arch} shows outperforms a more aggressive profile on this dataset size.

\subsection{Model Architecture}
\label{sec:arch}

The pipeline hosts four encoder families behind a shared interface. Every encoder emits a time-major sequence tensor that the CTC head of Section~\ref{sec:heads} consumes.

\subsubsection{CNN--RNN baseline}

The baseline is the CNN--BiLSTM--CTC network of Retsinas~\etal~\cite{retsinas2022htrbestpractices}: a ResNet-like CNN, column-wise max-pooling to collapse the height axis to a $128$-timestep sequence, and a three-layer BiLSTM (hidden $= 256$). Total parameters: $7.36$M.

\subsubsection{ViT-RGTS (main model)}

The main model is a CNN-stem Vision Transformer with register tokens. A four-layer convolutional stem collapses the $128 \times 1024$ input to a feature map of shape $D \times 1 \times 128$ with $D = 256$; the $128$ column descriptors form the patch-token sequence $p_1, \ldots, p_{128}$. A learnable tensor of $R$ \emph{register tokens} is prepended to this sequence:
\begin{equation}
\label{eq:token_sequence}
\mathbf{X}^{(0)} = [\mathbf{r}_1, \ldots, \mathbf{r}_R, \mathbf{p}_1, \ldots, \mathbf{p}_{128}] + \mathbf{E}_{\mathrm{pos}}[:S],
\end{equation}
where $S = R + 128$ and $\mathbf{E}_{\mathrm{pos}}$ is a learnable positional-embedding table. The register tokens are initialized from $\mathcal{N}(0, 0.02^2)$, matching Darcet~\etal~\cite{darcet2024registers}. When $R = 0$ only the patch tokens enter the encoder.

The encoder is a stack of $L = 6$ standard Transformer blocks~\cite{vaswani2017attention} with $h = 8$ heads, embedding width $D = 256$, feed-forward width $D_{\mathrm{ff}} = 1024$, and dropout $0.1$. The output patch tokens are passed to the CTC head; the register outputs are stored separately for the attention analysis. A parallel forward path exposes the per-layer attention tensors, per-token norms, and register embeddings needed by the interpretability metrics of Section~\ref{sec:attn_eval}. Total parameters: $9.48$M for $R = 4$; each additional register adds $D = 256$ parameters.

The CNN stem is a structural requirement rather than a refinement. It guarantees a strictly left-to-right token order along the width axis, which the CTC monotonic-alignment assumption depends on. Section~\ref{sec:res_abl_arch} shows that removing the stem (falling back to a row-major $16 \times 16$ patch grid) collapses CER by more than an order of magnitude.

\subsubsection{Pretrained comparators}

The first pretrained comparator is TorchVision's ImageNet-pretrained \texttt{vit\_b\_16}~\cite{dosovitskiy2021vit}, adapted to grayscale HTR input via a learnable $1{\times}1$ channel-expansion convolution. Fine-tuning uses a two-group differential learning rate (backbone $2 \times 10^{-5}$, head $5 \times 10^{-4}$). Total parameters: $91.9$M.

The second pretrained comparator is TrOCR-base~\cite{li2023trocr}, whose encoder is pretrained on real handwriting. The decoder is discarded; the encoder feeds a CTC head via aspect-ratio-preserving resize and a lightweight height-pooling adaptor. Fine-tuning uses layer-wise learning-rate decay~\cite{howard2018ulmfit}. Total parameters: $90.4$M.

For the READ2016 comparison a run of HTR-VT~\cite{li2024htrvt} is included as an independently designed CTC ViT external comparator. Total parameters: $8.75$M.

\subsubsection{CTC head}
\label{sec:heads}

All encoders share a dual-branch CTC head. A three-layer BiLSTM (hidden $= 256$) followed by a linear projection produces the primary character posteriors; a parallel $1{\times}3$ convolution with linear projection produces a shortcut posterior. Both are supervised by CTC loss~\cite{graves2006ctc},
\begin{equation}
\label{eq:ctc_loss}
\mathcal{L} = \mathcal{L}_{\mathrm{CTC}}(\log p_{\mathrm{rnn}}) + \lambda \, \mathcal{L}_{\mathrm{CTC}}(\log p_{\mathrm{cnn}}),
\end{equation}
with $\lambda = 0.1$; at inference only the BiLSTM branch is decoded via greedy best-path collapse. The dual branch follows the auxiliary shortcut of Retsinas~\etal~\cite{retsinas2022htrbestpractices} and is used across all four encoder families.

\subsubsection{Pipeline overview}

Figure~\ref{fig:pipeline} summarizes the end-to-end pipeline: a line image passes through preprocessing, one of the four encoder families, and the CTC head; the same forward pass exposes register tokens and attention tensors for interpretability analysis.

\begin{figure}[h]
\centering
\includegraphics[width=0.5\linewidth]{htr_pipeline_diagram}
\caption{End-to-end HTR pipeline. Register tokens (highlighted) are prepended to the patch-token sequence and shift the attention behaviour of the encoder without adding a decoder-side query.}
\label{fig:pipeline}
\end{figure}

\subsection{Training and Fine-Tuning}
\label{sec:training}

From-scratch models (CNN--RNN, ViT-RGTS, HTR-VT) use AdamW~\cite{loshchilov2019adamw} with weight decay $5 \times 10^{-4}$ and a differential learning-rate policy: CNN stem and CTC head at $1 \times 10^{-3}$, Transformer body at $5 \times 10^{-4}$. A five-epoch linear warmup is followed by cosine annealing to $10^{-6}$. All runs use $150$ epochs, seed $42$, deterministic cuDNN, effective batch $32$ (physical $8$, accumulation $\times 4$), and gradient clipping at $\ell_2$-norm $5$. The best-validation-CER epoch is retained for evaluation.

Pretrained models require a different treatment because a naive uniform learning rate destabilizes the pretrained weights (Section~\ref{sec:res_finetune}). Three strategies are supported: (i)~two-group differential LR --- uniform low LR on the backbone, higher LR on the head, used for ImageNet-pretrained ViT-B/16; (ii)~layer-wise learning-rate decay~\cite{howard2018ulmfit} --- $\eta_\ell = \eta_{\mathrm{base}} \cdot \gamma^{L-\ell}$ with $\gamma = 0.9$, used for TrOCR; and (iii)~gradual unfreezing, included for completeness and shown to fail for both models at the epoch budgets tested.

A stochastic weight-averaging (SWA) variant~\cite{izmailov2018swa} is trained in parallel for every family: from epoch $110$ of $150$ an averaged shadow model accumulates a running mean of the weights while the live network's learning rate is held constant via SWALR. Because the SWA-averaged model's evaluation is done after a separate BatchNorm-recalibration pass, its CER differs from the CER of the live network at the same epoch --- both numbers are collected but only the SWA-averaged evaluation is reported for SWA-enabled runs.

\subsection{Implementation and Hardware}
\label{sec:impl}

The pipeline is implemented in PyTorch~2.1 with mixed-precision AMP for the ViT models. All experiments ran on the NHR@FAU cluster (RTX~3080 partition, $10$~GB VRAM, $8$~CPU cores). A single $150$-epoch IAM run takes roughly $5$--$7$~hours for CNN--RNN, $9$--$11$~hours for ViT-RGTS, and $20$--$28$~hours for the fine-tuned pretrained lanes. Attention weights, register embeddings, and per-token norms are saved only for the best-validation checkpoint per run to keep storage bounded.
